This section is based on \cite{deift2006universalitymathematicalphysicalsystems}. 

Consider the scattering of neutrons of a heavy nucleus, say uranium $U^{238}$. The scattering cross-section is plotted as a function of the energy $E$ of the incoming neutrons, and one obtains a jagged graph with many hundreds of sharp peaks $E-1 < E_2 < \cdots$  and valleys $E'_1 < E'_2 < \cdots $. If $E \approx E_j$ for some $j$, the neutron is strongly repelled from the nucleus, and if $E \approx E'_j$ for some $j$, then the neutron sails through the nucleus, essentially unimpeded. The $E_j$’s are called scattering resonances. The challenge faced by physicists in the late $40$’s and early $50$’s was to develop an effective model to describe these resonances. One could of course write down a Schrodinger-type equation for the scattering system, but because of the high dimensionality of the problem there is clearly no hope of solving the equation for the $E_j$’s either analytically or numerically. However, as more experiments were done on heavy nuclei, each with hundreds of $E_j$’s, a consensus began to emerge that the “correct” theory of resonances was statistical, and here Wigner led the way. Any effective theory would have to incorporate two essential features present in the data, i.e.
\begin{enumerate}
	\item Modulo certain natural symmetry considerations, all nuclei in the same symmetry class exhibited universal behavior;
	\item In all cases, the $E_j$’s exhibited repulsion, or, more precisely, the probability that two $E_j$’s would be close together was small. 
\end{enumerate}

\begin{sic}
	(Eugene Wigner - On the Wigner Surmise (1956))
	“Perhaps I am now too courageous when I try to guess the distribution of the distances between successive levels (of energies of heavy nuclei). Theoretically, the situation is quite simple if one attacks the problem in a simpleminded fashion. The question is simply what are the distances of the characteristic values of a symmetric matrix with random coefficients.” 
\end{sic}

At some point in the mid-$1950$’s, in a striking development, Wigner suggested that the statistics of the neutron scattering resonances was governed by GOE (This is for time-reversal-invariant data, if this was not the case, this should be governed by the GUE). If one equips the space of real, symmetric matrices with a probability measure with a smooth density, the probability of a matrix $\m{M}$ having equal eigenvalues would be zero and the eigenvalues $\lambda_1, \cdots , \lambda_n$ of $\m{M}$ would comprise a random set with repulsion built in. So Wigner had a model, or more precisely, a class of models, which satisfied constraint $(ii)$. But why choose GOE? This is where the universality constraint $(i)$ comes into play. We quote from:
\begin{citation}
	(Wigner)
	“Let me say only one more word. It is very likely that the curve in Figure 1 is a universal function. In other words, it doesn’t depend on the details of the model with which you are working. There is one particular model in which the probability of the energy levels can be written down exactly. I mentioned this distribution already in Gatlinburg. It is called the Wishart distribution. Consider a set... .” 
\end{citation}
So in this way Wigner introduced GOE into theoretical physics: It provided a model with repulsion (and time-reversal) built in. Furthermore, the energy level distribution could be computed explicitly. By universality, it should do the trick! 